\documentclass[a4paper]{article}
\usepackage[warn]{mathtext}
\usepackage{graphicx} % Required for inserting images
\usepackage{wrapfig}
\usepackage{enumitem}
\usepackage[center]{caption}
\usepackage[english, russian]{babel}
\usepackage{amsmath}
\usepackage{geometry}
\setcounter{totalnumber}{5}

\geometry{verbose, a4paper, tmargin=2cm, bmargin=2cm, lmargin=2.0cm, rmargin=2.0cm}

\begin{document}
	
	\begin{titlepage}
		\centering
		{\scshape\Large МОСКОВСКИЙ ФИЗИКО-ТЕХНИЧЕСКИЙ ИНСТИТУТ \\
			(НАЦИОНАЛЬНЫЙ ИССЛЕДОВАТЕЛЬСКИЙ УНИВЕРСИТЕТ)\\ % название ВУЗа большим шрифтом
			Физтех-школа аэрокосмических технологий}
		

		\vspace{4cm} % отступ 4 см по вертикали
		{\LARGE Отчет о выполнении лабораторной работы 1.2.5}

		\vspace{1cm} % ещё 1 см
	{\huge\bf Измерение вынужденной регулярной \\прецессии гироскопа}
		
		\vspace{1cm} % ещё 1 см
		\vfill % заполнение по вертикали пустотой (LaTeX сам решит, сколько здесь отступить)
		
		\begin{flushright} % Этот текст будет с правого края страницы
			{\LARGE Губарев Никита, Б03-502}
		\end{flushright}
		\vfill 
		\today % Дата сборки документа
	\end{titlepage}
\tableofcontents 
\section{Аннотация}
Цели работы: исследовать вынужденную прецессию гироскопа; установить зависимость скорости вынужденной прецессии от величины момента сил, действующих на ось гироскопа; определить скорость вращения ротора гироскопа и сравнить её со скоростью, рассчитанной по скорости прецессии.

\section{Теоретические сведения}

Уравнения движения твердого тела можно представить в виде:
\begin{equation}
\frac{d\vec{P}}{dt} = \vec{F},
\end{equation}
\begin{equation}
\frac{d\vec{L}}{dt} = \vec{M}.
\end{equation} \\
Здесь (1) выражает закон движения центра масс тела, а (2) – уравнение моментов. Если сила $\vec{F}$ не зависит от угловой скорости, а момент сил $\vec{M}$ -- от скорости поступательного движения, то уравнения можно рассматривать независимо друг от друга. В данной работе для описания движения гироскопа потребуется только уравнение (2). \vspace{0.25 cm} \\
Момент импульса твёрдого тела в его главных осях можно выразить по формуле
\begin{equation}
\vec{L} = \vec{i}I_x\omega_x + \vec{y}I_y\omega_y + \vec{k}I_z\omega_z,
\end{equation} 
где $I_x, I_y, I_z$ – главные моменты инерции, а $\omega_x, \omega_y, \omega_z$ – компоненты вектора угловой скорости $\vec{\omega}$. \vspace{0.25 cm} \\ 
Быстро вращающееся тело, для которого произведение момента инерции на компоненту угловой скорости одной из осей много больше произведений для двух других осей, принято называть гироскопом.  \vspace{0.25 cm} \\
В силу (2) приращение момента импульса определяется интегралом
\begin{equation} 
\Delta \vec{L} = \int\vec{M} dt.
\end{equation}
Если момент внешних сил действует в течение короткого промежутка времени, из интеграла (4) следует, что приращение $\Delta\vec{L}$ момента импульса значительно меньще самого момента импульса. С этим связана устойчивость гироскопа, привденного в быстрое вращение. \vspace{0.25 cm}


\noindent Рассмотрим маховик, вращающийся вокруг оси $z$, перпендикулярной к его плоскости. Будем считать, что
\[
\begin{gathered}
\omega_z = \omega_0, \text{ } \omega_y = 0, \text { } \omega_x = 0.
\end{gathered}
\]
Пусть ось вращения повернулась в плоскости $zx$ по направлению к оси $x$ на бесконечно малый угол $d\varphi$. Такой поворот означает добавочное вращение маховика вокруг оси $y$, так что 
\[
\begin{gathered}
d\varphi = \Omega dt,
\end{gathered}
\]
где $\Omega$ -- угловая скорость вращения. Если $L_\Omega \ll L_{\omega_0}$ (5.1), то момент импульса маховика лишь повернется в плоскости $zx$ по направлению к оси $x$, не изменяя своей величины. Тогда изменение направлено вдоль оси $x$, и вектор $\vec{L}$ можно представить в виде векторного произведения вектора угловой скорости $\vec{\Omega}$ на вектор собственного момента испульса маховика $\vec{L}$. Таким образом, 
\begin{equation}
\frac{d\vec{L}}{dt} = \vec{\Omega} \times \vec{L}.
\end{equation}
Окончательно, в силу (2)
\begin{equation}
\vec{M} = \vec{\Omega} \times \vec{L}.
\end{equation}

\noindent Под действием момента $\vec{M}$ внешних сил ось гироскопа медленно вращается вокруг оси $y$ с угловой скоростью $\Omega$. Такое движение называется регулярной прецессией гироскопа. В данной работе для ее изучения на ось вращения гироскопа будут подвешены дополнительные грузы. Прецессию в таком случае вызывает момент силы тяжести (Приложение 1), а скорость ее равна
\begin{equation}
\Omega = \frac{mgl}{I_z\omega_0},
\end{equation}
где $m$ -- масса груза, $l$ -- расстояние от центра карданового подвеса до точки крепления груза. \vspace{0.25 cm} \\
Вычисление момента инерции ротора $I_z$ относительно оси вращения гироскопа происходит при помощи крутильного маятника. Так как период крутильных колебаний $T_0$ зависит лишь от момента инерции тела $I_0$ и модуля кручения проволоки
\[
\begin{gathered}
T_0 = 2\pi \sqrt{\frac{I_0}{f}},
\end{gathered}
\]
момент инерции $I_z$ можно рассчитать как
\begin{equation}
I_z = I_\text{ц} \frac{T_z^2}{T_\text{ц}^2},
\end{equation}
где $I_\text{ц}$ -- момент инерции циллиндра известных массы и диаметра; $T_\text{ц}, T_z$ -- периоды крутильных колебаний.
\section{Методика измерений}
В приложении 2 представлена схема установки.
\begin{enumerate}
	\item Установить ось гироскопа в горизонтальное положение, поворачивая ее за рычаг С.
	\item Включить питание гироскопа и подождать, чтобы вращение ротора успело стабилизироваться.
	\item Убедиться в том, что ротор вращается достаточно быстро. При легком воздействии на рычаг С, он не должен менять своего положения в пространстве.
	\item Подвесить к рычагу С груз Г. Должна начаться прецессия гироскопа. Трение в оси приводит к тому, что рычаг начнет медленно опускаться.
	\item Отклонить рычаг на $5-6^{\circ}$ вверх от горизонтальной плоскости. Подвесить к нему груз Г и с помощью секундомера найти угловую скорость прецессии $\Omega$ (по числу оборотов и времени прецессии). Измерения продолжать до тех пор, пока рычаг С не опустится на $5-6^{\circ}$ ниже горизонтальной плоскости, сделав целое количество оборотов относительно вертикальной оси. Измерить скорость опускания рычага С. Повторить опыт не менее пяти раз. Усреднить результаты.
	\item Проделать серию экспериментов, описанных в пункте 5, при различных значениях момента $M$ силы $F$ относительно центра масс гироскопа (длина плеча $l$ указана на установке). Результаты опытов представить в виде графика зависимости $\Omega$ от $M$.
	\item Измерить момент инерции ротора относительно оси симметрии $I_z$: подвесить ротор, извлеченный из такого же гироскопа, к концу вертикально висящей проволоки так, чтобы ось симметрии гироскопа была вертикальна; измерить период крутильных колебаний; заменить ротор гироскопа цилиндром, для которого известны радиус и масса; определить период крутильных колебаний для него. Пользуясь формулой (8), вычислить момент инерции ротора гироскопа $I_z$.
	\item Оценить погрешности $I_z$ и $\Omega$.
	\item Рассчитать с помощью (7) частоту вращения ротора.
	\item Определить частоту вращения ротора по фигурам Лиссажу. На "Вход Y" \ осциллографа подать сигнал со втором обмотки статора, а на "Вход синхр." \ сигнал с выхода генератора. Для получения эллипса: нажать тумблер "Вход X" ; \ повернуть "Стабильность" \ в крайнее правое положение. Подобрать частоту генератора так, чтобы эллипс стал неподвижным. Получение на экране осциллографа неподвижного эллипса означает, что частота сигнала генератора равна частоте вращения ротора гироскопа
	\item Оценить погрешность полученных результатов. Сравнить угловые скорости вращения ротора гироскопа, полученные разными методами.
	\item Убедиться в применимости соотношения (5.1)
\end{enumerate}
\section{Результаты измерений и обработка данных}

\subsection{Скорость вынужденной прецессии гироскопа}

\begin{gathered}
\overline{T} = \frac{1}{N}  \sum_{i=1}^N T_i ;  \text{ }\Omega = \frac{2\pi}{\overline{T}}
\end{gathered}


\begin{table}[!ht]
    \centering
    \begin{tabular}{|l|l|l|l|l|l|l|l|l|l|}
    \hline
        m, г & $М, Н \cdot м \cdot 10^-3$ & $\sigma_M$ & $Т_1, c$ &  $Т_2, c$ &  $Т_3, c$ &  $Т_{ср}, c$ & \sigma_{T_{ср}} & $\Omega, рад/с\cdot 10^-3$ & $\sigma_{\Omega}$  \\ \hline
        74,2 $\pm 0.1$& 88,0 & 0,1 & 140,5 & 140 & 140,5 & 140,3 & 3 & 45 & 1  \\ \hline
        92,4  $\pm 0.1$& 109,6 & 0,1 & 112 & 112 & 112,5 & 112,2 & 3 & 56 & 1  \\ \hline
        115,8  $\pm 0.1$& 137,3 & 0,1 & 90 & 89,5 & 89 & 89,5 & 3 & 70 & 2  \\ \hline
        138,0  $\pm 0.1$& 163,6 & 0,1 & 60,4 & 74,24 & 75 & 69,9 & 3 & 90 & 4  \\ \hline
        173,8  $\pm 0.1$& 206,1 & 0,1 & 59,7 & 59,7 & 59,7 & 59,7 & 3 & 105 & 5  \\ \hline
        213,3  $\pm 0.1$& 252,9 & 0,1 & 48,5 & 48,75 & 48,5 & 48,6 & 3 & 129 & 8  \\ \hline
    \end{tabular}
\end{table}

\noindent Т. к. зависимость $\Omega$ от $M = mgl$, где $l = 121$ мм, линейна, по методу наименьших квадратов,
Тогда $\frac{1}{I_z\omega_0} = 0,509 \pm 0,001$ с/кг. (Приложение 3)


\subsection{Момент инерции ротора}

Период крутильных колебаний ротора на проволоке составил $T_0 = 3,2$ с. Для цилиндра $T_\text{ц} = 4,0$ с. Радиус цилиндра $R_\text{ц} = 39,0$ мм, масса -- $m_\text{ц} = 1616,0$ г. Тогда его момент инерции:
\[
\begin{gathered}
I_\text{ц} = \frac{m_\text{ц} r^2}{2} \approx 1,229 \cdot 10^{-3} \text{ кг} \cdot \text{м}^2.
\end{gathered}
\]
Из соотношения (8) момент инерции ротора:
\[
\begin{gathered}
I_z = I_\text{ц} \frac{T_0^2}{T_\text{ц}^2} \approx 0,787  \cdot 10^{-3} \text{ кг} \cdot \text{м}^2.
\end{gathered}
\]
Погрешности вычисления моментов инерции:
\[
\begin{gathered}
\sigma_{I_{\text{ц}}} = \frac{1}{2} I_{\text{ц}} \sqrt{2(\frac{\Delta_{\text{шт}}}{R})^2 + (\frac{\Delta_{\text{в}}}{m_{\text{ц}}})^2} \approx 0,001 \cdot 10^{-3}  \text{ кг} \cdot \text{м}^2,
\end{gathered}
\]
\[
\begin{gathered}
\sigma_{I_{z}} = I_{z} \sqrt{2(\frac{\Delta_{c}}{T_{\mathrm{ц}}})^{2} + 2(\frac{\Delta_{c}}{T_{0}})^{2} + (\frac{\sigma_{I_{\mathrm{ц}}}}{I_{\mathrm{ц}}})^{2}} \approx 0,003 \cdot 10^{-3}  \text{ кг} \cdot \text{м}^2.
\end{gathered}
\]

\subsection{Угловая скорость вращения ротора гироскопа}

\[
\begin{gathered}
\omega_0 = \frac{1}{kI_z} \approx 2496 \text{ рад/c}; \text{ } \sigma_{\omega_{0}} = \omega_{0} \sqrt{(\frac{\sigma_{k}}{k})^{2} + (\frac{\sigma_{I_{z}}}{I_{z}})^{2}} = 11 \text{ рад/с}.
\end{gathered}
\]
Частота вращения ротора гироскопа также измерялась осциллографом при помощи фигур Лиссажу: $\nu \approx 400$ Гц. Отсюда $\omega = 2 \pi \nu \approx 2513$ рад/с.

\section{Выводы}

Полученные разными методами значения угловой скорости вращения ротора совпали в пределах погрешности:

Угловая скорость, полученная при помощи измерения прецессии: $2496 рад/с$

Угловая скорость, полученная при помощи измерения резонансной частоты от генератора сигналов: $2513 рад/с$

При этом относительная погрешность полученного результата невелика:  $\varepsilon \approx 0,7 $\%. Все теоретические закономерности, рассмотренные в данной работе, также выполнились в пределах точности эксперимента.


\begin{wrapfigure}{R}{0.3\textwidth}
    \includegraphics[width=1\linewidth]{"C:/Users/yaros/OneDrive/Desktop/repositories/mipt_nik/Физос/Лабы/1.2.5/1.2.5_scheme.png}
    \caption{Приложение 1. Иллюстрация направления сил в эксперименте}
\end{wrapfigure}
\section{Приложения}
\begin{wrapfigure}{R}{0.3\textwidth}
    \includegraphics[width=1\linewidth]{"C:/Users/yaros/OneDrive/Desktop/repositories/mipt_nik/Физос/Лабы/1.2.5/1.2.5_gyro.png}
    \caption{Приложение 2. Схема маховика}
\end{wrapfigure}


\begin{figure}[!h]
   \centering    
    \includegraphics[width=1\linewidth]{C:/Users/yaros/OneDrive/Desktop/repositories/mipt_nik/Физос/Лабы/1.2.5/Figure_1.png}
    \caption{Приложение 3. График зависимости угловой скорости прецессии от приложенного момента силы тяжести} 
\end{figure}
\end{document}




